\chapter{Building the Documentation}
\label{chap:building-docs}

\lupnt{} ships two documents built from the same repository: the Sphinx
website under \file{docs/}, which combines narrative pages, the executed
tutorial notebooks, the Python API, and a Doxygen-derived C++ API; and this
LaTeX manual under \file{latex_docs/}, which carries the mathematical
specifications in full. This chapter specifies how each is produced, which
files configure it, what the build writes where, and how to extend either one.
It expands \file{docs/builddocs.rst} with the contents of
\file{docs/make\_docs.py}, \file{docs/conf.py}, \file{docs/index.rst},
\file{docs/CMakeLists.txt}, \file{docs/Doxyfile}, and the LaTeX sources
themselves.

The developer environment that supplies the toolchain --- Pixi, its tasks, and
the CI pipeline that publishes the site --- is specified in
\cref{chap:development}.

\section{The documentation toolchain}
\label{sec:bd-toolchain}

The narrative documentation and the Python API pages are written in
reStructuredText and rendered by Sphinx. The tutorial pages come from two
sources: the Python example notebooks, included through \code{nbsphinx} and
\code{nbsphinx-link}, and the C++ example programs, included with
\code{literalinclude}. The C++ API reference is generated from Doxygen
comments in the public headers and pulled into Sphinx by the Breathe and Exhale
extensions. The generated HTML is GitHub Pages-ready.

\begin{table}[H]
  \centering
  \caption{Documentation toolchain components and where they are configured.}
  \label{tab:bd-toolchain}
  \small
  \begin{tabularx}{\textwidth}{@{}l L L@{}}
    \toprule
    Component & Role & Configured in \\
    \midrule
    Sphinx           & site generator, HTML builder & \file{docs/conf.py} \\
    \code{furo}      & HTML theme (light and dark logos) & \file{docs/conf.py} \\
    \code{autodoc}, \code{autosummary}, \code{napoleon}
                     & Python API extraction, Google/NumPy docstrings & \file{docs/conf.py} \\
    \code{nbsphinx}, \code{nbsphinx-link}
                     & renders \file{.ipynb} tutorials from stored outputs
                     & \file{docs/conf.py}, \file{docs/tutorial/Python/*.nblink} \\
    Doxygen          & parses the C++ headers into XML
                     & \code{exhaleDoxygenStdin} in \file{docs/conf.py} \\
    \code{breathe}   & bridges Doxygen XML into the Sphinx C++ domain & \file{docs/conf.py} \\
    \code{exhale}    & generates the C++ API page tree & \file{docs/conf.py} \\
    \code{pandoc}    & notebook markdown conversion for \code{nbsphinx} & Pixi dependency \\
    \code{tectonic}  & builds this LaTeX manual & \file{latex_docs/build.sh} \\
    \bottomrule
  \end{tabularx}
\end{table}

The layout of \file{docs/} is:

\begin{lstlisting}[language=shell]
docs/
  index.rst              # root toctree, and the MAKE_DOCS python-api markers
  conf.py                # Sphinx configuration (incl. the Doxygen stdin config)
  make_docs.py           # the build driver invoked by `pixi run build-docs`
  make.bat               # Windows convenience wrapper around sphinx-build
  requirements.txt       # pip requirements for a non-pixi build
  Doxyfile               # standalone Doxygen config, used by docs/CMakeLists.txt
  CMakeLists.txt         # the GenerateDocs (m.css) target
  .nojekyll              # copied into the published site
  introduction.rst  getting_started.rst  development.rst  builddocs.rst
  initialization.rst  roadmap.md
  pages/                 # new_simulation, sp3_download, cross_validation,
                         #   related_publications, about.dox
  math/                  # the mathematical specification pages (index + 11 pages)
  tutorial/
    Python/              # one .nblink per python/examples/exN_*.ipynb
    C++/                 # one .rst per cpp/examples tutorial program
  _static/               # logos, background images, generated cesium scenes
  latex/                 # THIS manual (see the last section of this chapter)
\end{lstlisting}

\section{Building with Pixi}
\label{sec:bd-pixi}

One task builds everything:

\begin{lstlisting}[language=shell]
pixi run build-docs
\end{lstlisting}

The task is defined as \code{python docs/make\_docs.py} with
\code{depends-on = ["build-py"]}. The dependency matters: the Python API
generator \emph{imports} \pylupnt{} to enumerate its classes and functions, so
the compiled extension must exist. If \file{python/pylupnt/\_pylupnt*.so} is
missing --- which only happens on a manual, non-Pixi build --- build the
bindings first:

\begin{lstlisting}[language=shell]
pixi run build-py
\end{lstlisting}

\subsection{What \file{make\_docs.py} does}
\label{sec:bd-makedocs}

\file{docs/make\_docs.py} descends from the Open3D documentation driver,
modified to use Breathe and Exhale for the C++ API. Running it performs the
following steps, in order, with the working directory set to \file{docs/}:

\begin{enumerate}[leftmargin=1.6em]
  \item Create or clear the staging directory \file{docs/\_out}.
  \item Remove \file{docs/\_static/C++/build} if present.
  \item With \code{--clean} (the default), remove \file{docs/\_out},
    \file{docs/cpp\_api}, \file{docs/python\_api}, and \file{docs/\_build}.
  \item Read \file{docs/index.rst} and collect the modules to document from the
    \code{MAKE\_DOCS/python\_api/...} markers (\cref{sec:bd-pythonapi}).
  \item For every such module, generate one \file{.rst} file per class, per
    function, and one summary page per submodule, into \file{docs/python\_api}.
  \item Invoke \code{sphinx-build -b html -j <jobs> . \_out/html}. Exhale runs
    Doxygen from inside this call and generates \file{docs/cpp\_api}.
  \item With \code{--copy} (the default), delete \file{build/docs}, copy
    \file{docs/\_out/html} there, remove the \file{.doctrees} directory, and
    copy \file{docs/.nojekyll} into it.
  \item Clean the staging directories again.
\end{enumerate}

The final site is therefore the repository-level \file{build/docs} directory,
not anything under \file{docs/}. The \file{.nojekyll} marker is what makes
GitHub Pages serve generated directories whose names begin with an underscore,
such as \file{\_static}, instead of filtering them through Jekyll.

Sphinx parallelism defaults to \code{multiprocessing.cpu\_count()} and is
overridden by the environment variable \code{LUPNT\_DOCS\_JOBS}:

\begin{lstlisting}[language=shell]
LUPNT_DOCS_JOBS=4 pixi run build-docs
\end{lstlisting}

The CI \code{pages} job sets exactly this variable to \code{4}
(\cref{sec:bd-ci}).

\begin{lupntnote}
\file{make\_docs.py} hard-codes the \code{html} builder. To produce any other
output --- \code{latex}, \code{linkcheck}, \code{dirhtml} --- call
\code{sphinx-build} directly from \file{docs/}, after generating the Python API
pages at least once so the \code{python\_api/*} toctree entries resolve.
\end{lupntnote}

\subsection{Sphinx configuration}
\label{sec:bd-conf}

\file{docs/conf.py} inserts \file{<repo>/python} at the front of
\code{sys.path} so \code{import pylupnt} resolves from the source tree exactly
as it does under Pixi activation, then declares:

\begin{lstlisting}[language=py]
# docs/conf.py
extensions = [
    "sphinx.ext.autodoc",
    "sphinx.ext.autosummary",
    "sphinx.ext.napoleon",
    "breathe",
    "exhale",
    "nbsphinx",
    "nbsphinx_link",
]

html_theme = "furo"
html_title = "LuPNT"
html_static_path = ["_static"]
html_favicon = "_static/lupnt_mark.svg"
html_theme_options = {
    "light_logo": "lupnt_logo_horizontal.svg",
    "dark_logo": "lupnt_logo_horizontal_dark.svg",
}

autodoc_default_options = {
    "members": True,
    "undoc-members": True,
    "inherited-members": True,
}
autosummary_generate = True
napoleon_google_docstring = True
napoleon_numpy_docstring = True
nbsphinx_execute = "never"

primary_domain = "cpp"
highlight_language = "cpp"
\end{lstlisting}

Because \code{primary\_domain} and \code{highlight\_language} are both
\code{cpp}, an unqualified code role or code block in a \file{.rst} page is
interpreted and highlighted as C++; Python snippets must say so explicitly.

\subsection{The C++ API: Breathe and Exhale}
\label{sec:bd-cppapi}

Exhale, not the committed \file{docs/Doxyfile}, drives the Doxygen run used by
the website. \code{exhaleUseDoxyfile} is \code{False} and
\code{exhaleExecutesDoxygen} is \code{True}, so the configuration is supplied
inline:

\begin{lstlisting}[language=py]
# docs/conf.py
breathe_projects = {"lupnt": str(DOCS_DIR / "_build" / "doxygen" / "xml")}
breathe_default_project = "lupnt"

exhale_args = {
    "containmentFolder": "./cpp_api",
    "rootFileName": "cpp_library_root.rst",
    "rootFileTitle": "C++ API",
    "doxygenStripFromPath": str(REPO_ROOT / "cpp" / "lupnt"),
    "createTreeView": True,
    "exhaleExecutesDoxygen": True,
    "exhaleUseDoxyfile": False,
    "exhaleDoxygenStdin": """
PROJECT_NAME = LuPNT
INPUT = <repo>/cpp/lupnt
RECURSIVE = YES
FILE_PATTERNS = *.h
EXCLUDE_PATTERNS = */environment/plasma/fortran/*
GENERATE_XML = YES
GENERATE_HTML = NO
GENERATE_LATEX = NO
XML_OUTPUT = xml
EXTRACT_ALL = YES
EXTRACT_ANON_NSPACES = NO
HIDE_UNDOC_MEMBERS = NO
HIDE_UNDOC_CLASSES = NO
QUIET = YES
""",
}
\end{lstlisting}

Two of these settings encode deliberate decisions:

\begin{itemize}[leftmargin=1.4em]
  \item \code{FILE\_PATTERNS = *.h} with \code{EXTRACT\_ANON\_NSPACES = NO}
    documents only the public headers. Implementation files hold internal
    anonymous-namespace helpers that Doxygen would otherwise surface as
    \code{lupnt::@<number>} namespaces in the API listing; the public API lives
    in the headers.
  \item \code{EXCLUDE\_PATTERNS} drops
    \file{cpp/lupnt/environment/plasma/fortran}, whose legacy fixed-form
    Fortran sources are not meaningfully documentable by Doxygen.
\end{itemize}

The generated pages land in \file{docs/cpp\_api}, rooted at
\file{cpp\_api/cpp\_library\_root.rst}, and the Doxygen XML in
\file{docs/\_build/doxygen/xml}. Both are regenerated from scratch on every
build and are removed by the clean steps of \cref{sec:bd-makedocs}, so neither
is committed.

To document a new C++ symbol on the website, no documentation file is needed:
add a Doxygen comment to its declaration in a header under \file{cpp/lupnt},
and it appears in the tree on the next build.

\subsection{The Python API pages}
\label{sec:bd-pythonapi}

Which Python modules get API pages is declared in a comment block inside
\file{docs/index.rst}. \file{make\_docs.py} parses lines matching
\code{MAKE\_DOCS/python\_api/<module> <type>}; the optional second field is
\code{python\_only} for pure-Python modules, and is omitted for
pybind11-backed ones. The distinction changes the member-discovery rule:
\code{python\_only} modules are scanned for classes and functions whose
\code{\_\_module\_\_} equals the module name, while binding modules are scanned
for classes belonging to the module or a submodule and for built-in functions.

\begin{lstlisting}[language=py]
# docs/index.rst (inside a comment block)
MAKE_DOCS/python_api/pylupnt
MAKE_DOCS/python_api/pylupnt.plot python_only
MAKE_DOCS/python_api/pylupnt.core.pylupnt_utils python_only
MAKE_DOCS/python_api/pylupnt.plasma.kp_loader python_only
\end{lstlisting}

The same four modules are listed again, without the marker prefix, in the
\emph{Python API} toctree so that the generated pages are actually reachable:

\begin{lstlisting}[language=py]
# docs/index.rst
.. toctree::
    :maxdepth: 1
    :caption: Python API

    python_api/pylupnt
    python_api/pylupnt.plot
    python_api/pylupnt.core.pylupnt_utils
    python_api/pylupnt.plasma.kp_loader
\end{lstlisting}

Adding a module to the Python API therefore requires editing both places: the
marker block (so the pages are generated) and the toctree (so they are linked).
Forgetting the first yields a broken toctree entry; forgetting the second
yields orphaned pages and a Sphinx warning.

\subsection{The site toctrees}

\file{docs/index.rst} groups the pages into four captioned toctrees, plus the
Python and C++ API ones above:

\begin{table}[H]
  \centering
  \caption{Top-level structure of \file{docs/index.rst}.}
  \label{tab:bd-toctrees}
  \small
  \begin{tabularx}{\textwidth}{lL}
    \toprule
    Caption & Entries \\
    \midrule
    Getting Started & \file{introduction}, \file{getting\_started},
      \file{development}, \file{pages/new\_simulation}, \file{builddocs},
      \file{pages/sp3\_download}, \file{pages/cross\_validation} \\
    Tutorial & \file{tutorial/Python/index}, \file{tutorial/C++/index},
      \file{pages/related\_publications} \\
    Math Specs & \file{math/index} \\
    Python API / C++ API & the generated trees of
      \cref{sec:bd-pythonapi,sec:bd-cppapi} \\
    \bottomrule
  \end{tabularx}
\end{table}

This is the same ordering the LaTeX manual follows in its four parts, so a page
added to one document has an obvious home in the other.

\section{Previewing locally}
\label{sec:bd-preview}

Serve the generated site rather than opening the HTML files from disk:

\begin{lstlisting}[language=shell]
python -m http.server 8000 --directory build/docs
\end{lstlisting}

then open \url{http://localhost:8000}. A local server is required for in-page
search and for the interactive visualizations to load; opening
\file{build/docs/index.html} directly leaves both broken because of the
browser's file-origin restrictions.

\section{Tutorial outputs}
\label{sec:bd-notebooks}

\code{nbsphinx} is configured with \code{nbsphinx\_execute = "never"}. The
rendered notebooks therefore display the outputs \emph{stored in} each
\file{.ipynb}; the documentation build never re-runs them. This is a
deliberate trade: the docs build stays fast and needs no runtime data bundle,
at the cost of the stored outputs being a committed artefact that can go stale.

After changing an example, or an API an example depends on, refresh those
stored outputs in place:

\begin{lstlisting}[language=shell]
pixi run render-notebooks                 # all notebooks
pixi run render-notebooks -- --only ex16  # a subset
pixi run render-tutorials                 # render, then optimize the PNGs
\end{lstlisting}

The \code{nbstripout} pre-commit hook is configured to preserve outputs for
exactly these files (\file{python/examples/exN\_*.ipynb}) and to strip every
other notebook in the repository; see \cref{sec:dev-style}.

Three classes of output behave differently once published:

\begin{itemize}[leftmargin=1.4em]
  \item \textbf{Static Matplotlib figures} are embedded as images and always
    render, offline included.
  \item \textbf{Plotly figures} (examples \code{ex12}, \code{ex14} through
    \code{ex16}) pull \code{plotly.js} from a CDN at view time, so a network
    connection is required to see them.
  \item \textbf{Cesium scenes} (example \code{ex13}) are written as
    self-contained HTML into \file{docs/\_static/cesium\_scenes} and linked
    from the page; they pull CesiumJS from a CDN.
\end{itemize}

The headless-Chrome prerequisite for the static Plotly renders, and the
system libraries it needs on a minimal Linux image, are documented in
\cref{sec:dev-notebooks}.

\section{Manual build without Pixi}
\label{sec:bd-manual}

If Pixi is not available, build the C++ library and the Python bindings first
(\cref{sec:dev-bindings}), then install the documentation dependencies and run
the driver:

\begin{lstlisting}[language=shell]
cd docs
pip install -r requirements.txt
brew install doxygen pandoc          # macOS
sudo apt-get install doxygen pandoc  # Linux
python make_docs.py
\end{lstlisting}

\file{docs/requirements.txt} lists \code{notebook}, \code{sphinx},
\code{sphinx-rtd-theme}, \code{furo}, \code{exhale}, \code{breathe},
\code{nbsphinx}, \code{nbsphinx-link}, \code{jinja2}, and
\code{lxml\_html\_clean}; Doxygen and pandoc are not pip-installable and must
come from a system package manager or conda-forge. As with the Pixi build, the
generated HTML is written to the repository-level \file{build/docs} folder.

\begin{lupntwarning}
A manual build must be able to \code{import pylupnt}. Either run it inside a
shell where \code{PYTHONPATH} contains \file{python/}, or install the package;
otherwise the Python API generation step of \cref{sec:bd-makedocs} raises an
import error before Sphinx is ever invoked.
\end{lupntwarning}

\section{The standalone Doxygen target}
\label{sec:bd-doxygen}

Independently of the Sphinx site, \file{docs/CMakeLists.txt} defines a
\code{lupnt\_docs} CMake project with a \code{GenerateDocs} target that renders
Doxygen through \code{m.css} for a standalone, more richly styled C++ reference:

\begin{lstlisting}[language=shell]
cmake -S docs -B build-docs-cmake -GNinja
cmake --build build-docs-cmake --target GenerateDocs
\end{lstlisting}

The project fetches \code{mosra/m.css} at a pinned commit through CPM, adds the
\file{cpp} project as a dependency, configures \file{docs/Doxyfile} and
\file{docs/conf.py} into the binary directory, and runs
\code{m.css/documentation/doxygen.py} on the configured \file{conf.py},
writing to \code{<binary dir>/doxygen}. The variables substituted into the
\file{Doxyfile} are \code{DOXYGEN\_PROJECT\_NAME},
\code{DOXYGEN\_PROJECT\_VERSION}, \code{DOXYGEN\_PROJECT\_ROOT}, and
\code{DOXYGEN\_OUTPUT\_DIRECTORY}.

\begin{lupntnote}
The committed \file{docs/Doxyfile} (project name \code{"LuPNT (C++ API)"},
\code{OUTPUT\_DIRECTORY = doxygen}) is used by this target only. The website's
C++ API uses the inline Doxygen configuration of \cref{sec:bd-cppapi} instead,
so changing the \file{Doxyfile} has no effect on \code{pixi run build-docs}.
\end{lupntnote}

\section{Publishing from CI}
\label{sec:bd-ci}

The \code{pages} job in \file{.github/workflows/} is the only job in the
\code{docs} stage and runs only on the default branch:

\begin{lstlisting}[language=yamlcfg]
# .github/workflows/
pages:
  extends: .build_template
  stage: docs
  variables:
    LUPNT_DOCS_JOBS: "4"
  script:
    - *download_lupnt_data
    - pixi run build-docs
    - rm -rf public
    - mv build/docs public
  artifacts:
    paths:
      - public
  rules:
    - if: '$CI_COMMIT_BRANCH == $CI_DEFAULT_BRANCH'
\end{lstlisting}

It downloads the runtime data bundle first, because \code{build-docs} depends
on \code{build-py} and the import of \pylupnt{} performed by the Python API
generator touches data paths. The site is then renamed to \file{public}, which
is where GitHub Pages expects it, and published as a job artifact. The
\file{.nojekyll} file that \file{make\_docs.py} copies makes the same tree
serve correctly from GitHub Pages.

\section{This manual: \file{latex_docs}}
\label{sec:bd-latex}

The LaTeX manual is a separate document with a separate build, and it lives at
the repository root rather than under \file{docs/}. That placement is
deliberate: the manual is a self-contained reference, so \file{docs/} may
eventually be retired, and the manual is designed to survive that. It reads nothing from \file{docs/} --- the
\code{\textbackslash graphicspath} entry for \file{docs/_static} is an optional
fallback, not a dependency.

It carries the mathematical specifications as full derivations, unit contracts,
and inline implementing code, together with the bibliography behind the models
--- of which the primary reference is \citep{iiyama2026thesis}. What only the
Sphinx site provides is the auto-generated API reference and the executed
notebooks; those are generated artefacts and are not duplicated here.

\subsection{Layout}

\begin{table}[H]
  \centering
  \caption{Files under \file{latex_docs}.}
  \label{tab:bd-latex-layout}
  \small
  \begin{tabularx}{\textwidth}{lL}
    \toprule
    Path & Role \\
    \midrule
    \file{lupnt.tex}    & the master document: class options, title page, front
      matter, the four \code{\textbackslash part} divisions, and one
      \code{\textbackslash include} per chapter \\
    \file{preamble.tex} & every package, listing style, TikZ style, admonition
      environment, and macro available to a chapter --- the only place
      \code{\textbackslash usepackage} may appear \\
    \file{chapters/}    & one \file{.tex} fragment per chapter, each
      \code{\textbackslash include}d by \file{lupnt.tex} \\
    \file{refs.bib}     & the bibliography database; the only source of
      citation keys \\
    \file{build.sh}     & the build driver; also resolves
      \code{LUPNT\_DATA\_PATH} and writes \file{figpath.tex} \\
    \file{STYLE.md}     & the chapter style contract every fragment must follow \\
    \file{build/}       & generated output, including \file{build/lupnt.pdf} \\
    \bottomrule
  \end{tabularx}
\end{table}

\subsection{Building}

No system LaTeX installation is required. \file{build.sh} uses \code{tectonic},
which downloads the packages it needs on demand and runs BibTeX itself, so a
single invocation produces a document with resolved citations and cross
references:

\begin{lstlisting}[language=shell]
cd latex_docs
./build.sh            # full build into ./build
./build.sh --keep     # keep .aux/.bbl intermediates for debugging
\end{lstlisting}

The script prefers a \code{tectonic} already on \code{PATH} and otherwise falls
back to \code{pixi exec}, which fetches it into a throwaway environment ---
nothing is added to \file{pixi.toml} or the lockfile:

\begin{lstlisting}[language=shell]
# latex_docs/build.sh, in effect
if command -v tectonic >/dev/null 2>&1; then
  RUNNER=(tectonic)
elif command -v pixi >/dev/null 2>&1; then
  RUNNER=(pixi exec --spec tectonic tectonic)
fi
"${RUNNER[@]}" --outdir build lupnt.tex
\end{lstlisting}

Without the \code{--keep} flag the script deletes the \code{.aux}, \code{.log},
\code{.out}, \code{.toc}, \code{.bbl}, \code{.blg}, and \code{.bcf}
intermediates, leaving only \file{build/lupnt.pdf}. Equivalent invocations, if
either tool is preferred directly:

\begin{lstlisting}[language=shell]
# tectonic, directly
tectonic --outdir build lupnt.tex

# or a classic TeX Live toolchain
pdflatex lupnt && bibtex lupnt && pdflatex lupnt && pdflatex lupnt
\end{lstlisting}

\begin{lupntnote}
\file{refs.bib} must sit beside \file{lupnt.tex}, which is why it is committed
in \file{latex_docs} rather than generated or symlinked from elsewhere:
tectonic resolves the bibliography relative to the main file.
\end{lupntnote}

The manual is not built by \code{pixi run build-docs} and is not published by
the \code{pages} job. It is built on demand.

\subsection{Adding a chapter}
\label{sec:bd-newchapter}

Three steps, in order:

\begin{enumerate}[leftmargin=1.6em]
  \item Write \file{latex_docs/chapters/foo.tex} as a \emph{fragment}. It must
    not contain \code{\textbackslash documentclass},
    \code{\textbackslash usepackage}, \code{\textbackslash begin\{document\}},
    or a bibliography --- \file{preamble.tex} has already been read by the time
    the fragment is included. It opens with a
    \code{\textbackslash chapter\{...\}} and its label:
\begin{lstlisting}
% latex_docs/chapters/foo.tex
\chapter{Time Systems and Conversions}
\label{chap:time-conversions}

Opening paragraph: what this chapter specifies, which source files
implement it, and which references it follows.

\section{...}
\label{sec:tc-...}
\end{lstlisting}
  \item Add \code{\textbackslash include\{chapters/foo\}} to
    \file{lupnt.tex}, inside the \code{\textbackslash part} it belongs to. The
    order of the \code{\textbackslash include} lines is the order of the
    document; there is no separate table of contents to maintain.
  \item Follow \file{STYLE.md}. It is the contract, not a suggestion; the
    points it fixes are summarised below.
\end{enumerate}

If a new package is genuinely needed, add it to \file{preamble.tex} rather than
to the chapter, so that every chapter stays compilable in isolation via
\code{\textbackslash includeonly}.

\subsection{The style contract in brief}

\file{STYLE.md} is authoritative; this is an index of what it fixes.

{\small
  \begin{longtable}{@{}>{\raggedright\arraybackslash}p{0.40\textwidth}
                      >{\raggedright\arraybackslash}p{0.54\textwidth}@{}}
    \caption{Macros and environments a chapter may use, from
             \file{preamble.tex}.}
    \label{tab:bd-style} \\
    \toprule
    Construct & Use for \\
    \midrule
    \endfirsthead
    \toprule
    Construct & Use for \\
    \midrule
    \endhead
    \bottomrule
    \endlastfoot
    \code{\textbackslash code\{x\}} & identifier, enum value, CLI flag, short
      expression \\
    \code{\textbackslash cpp\{X\}}, \code{\textbackslash py\{x\}} & a C++ or
      Python symbol \\
    \code{\textbackslash file\{cpp/lupnt/x.cc\}} & a repository path (breaks at
      \code{/} and the underscore) \\
    \code{\textbackslash lupnt\{\}}, \code{\textbackslash pylupnt\{\}} & the
      library name, the Python package \\
    \code{\textbackslash vec}, \code{\textbackslash mat},
      \code{\textbackslash uvec} & 3-vector, matrix, unit vector \\
    \code{\textbackslash dd}, \code{\textbackslash transpose} & upright
      differential, transpose superscript \\
    \code{\textbackslash diag}, \code{\textbackslash atantwo},
      \code{\textbackslash argmin}, \code{\textbackslash Tr},
      \code{\textbackslash erf} & operators \\
    \code{\textbackslash TT}, \code{\textbackslash TAI},
      \code{\textbackslash TDB}, \ldots & time scales, in math mode \\
    \code{\textbackslash SI}, \code{\textbackslash si},
      \code{\textbackslash num} & units, via \code{siunitx}; extra units
      \code{\textbackslash yr}, \code{\textbackslash dayunit},
      \code{\textbackslash bit}, \code{\textbackslash dB},
      \code{\textbackslash dBHz}, \code{\textbackslash dBW},
      \code{\textbackslash dBi}, \code{\textbackslash TECU},
      \code{\textbackslash au} \\
    \code{lupntnote}, \code{lupntwarning} & replace the reStructuredText
      \code{.. note::} and \code{.. warning::} \\
    \code{lstlisting} & languages \code{cpp}, \code{py}, \code{yamlcfg},
      \code{shell}, \code{bibtexlang} \\
    \code{booktabs}, \code{tabularx} with the ragged-right \code{L} column &
      tables \\
    TikZ styles \code{tsnode}, \code{boxnode} and edge classes
      \code{constant}, \code{table}, \code{linear}, \code{model},
      \code{integral}, \code{flow} & diagrams \\
  \end{longtable}
}

Four rules are easy to violate and worth restating:

\begin{itemize}[leftmargin=1.4em]
  \item \textbf{Labels are prefixed by kind and scope}: \code{chap:<slug>},
    \code{sec:<slug>-<name>}, \code{eq:<slug>-<name>},
    \code{tab:<slug>-<name>}, \code{fig:<slug>-<name>},
    \code{lst:<slug>-<name>}. Reference them with
    \code{\textbackslash cref} or \code{\textbackslash Cref}, never with a
    hard-coded number.
  \item \textbf{Citations use keys that already exist in \file{refs.bib}}. Add
    an entry rather than inventing a key. Chapter and equation numbers in the
    thesis are cited as
    \code{\textbackslash citep[Ch.~6.2]\{iiyama2026thesis\}}.
  \item \textbf{Listings are ASCII only} --- no arrows, Greek letters, or en
    dashes --- and every listing names the file it came from on its first
    line. The same applies to prose: non-ASCII characters belong in math mode,
    as \code{\$\textbackslash mu\$} or \code{\$\textbackslash approx\$}, never
    bare.
  \item \textbf{Chapters are specifications, not summaries.} Each model states
    the state, unit, and time-scale contract it assumes; the governing
    equations, numbered, with every symbol defined; the implementing code with
    its source path; the configuration knobs that change its behaviour; and an
    explicit \emph{model boundaries} section covering what is approximated,
    what is omitted, and what accuracy is expected.
\end{itemize}

\subsection{Keeping the two documents consistent}

The Sphinx site and this manual are written from the same repository but are
not generated from one another. Consequences worth planning around:

\begin{itemize}[leftmargin=1.4em]
  \item A change to a \file{.rst} page does not propagate to the corresponding
    chapter, and vice versa. The mapping is one chapter per page, with the
    filenames deliberately kept parallel
    (\file{docs/development.rst} $\leftrightarrow$
    \file{latex_docs/chapters/development.tex}).
  \item The manual has no CI job. A LaTeX error is caught only when someone
    runs \file{build.sh}; there is no equivalent of the \code{pages} job to
    fail the pipeline.
  \item Undefined cross references and missing citations are warnings, not
    errors, under tectonic. Check the build log for
    \code{LaTeX Warning: Reference ... undefined} and
    \code{Citation ... undefined} before considering a chapter finished.
  \item Notebook outputs and the API references appear only on the website. The
    manual draws its own figures in TikZ instead, so that they add no binary
    weight to the repository and cannot drift out of sync with the prose.
  \item \file{latex_docs} is \textbf{text only}; its \file{.gitignore} rejects
    images and PDFs. Assets shared with the Sphinx site are read from
    \file{docs/_static}, and generated or large artefacts from
    \file{\$LUPNT_DATA_PATH/docs}, included with \code{\textbackslash datafig}
    so that a clone without the fetched data directory still builds --- the
    figure degrades to a placeholder rather than failing.
\end{itemize}
